\scanpage{94}

J'ai découvert un canular ennuyeux dans EGA IV 17.1.6 (i), concernan le cas ``formellement lisse''. En effet, le démonstration via (16.5.17) ne marche que si on suppose $\quad{}^1_{X/Y}$ de présentation finie, par exemple si $f$ est localement de \deleted{présentati} type fini ! J'\deleted{\unsure{n|e}}ignore donc si la notion ``formellement lisse'' est vraiment locale en haut ! On peut montrer que la question équivaut à la suivante: si $A$ est un anneau commutatif, et $M$ un $A$-Module (pas nécessairement de présentation finie) qui est ``localement projectif sur $\operatorname{Spec}(A)$'', i.e. tel qu'il existe des $f_i\quad A$ engendrant l'idéal unité, avec $M_{f_i}$ projectif sur $A_{f_i}$ , alors $M$ est-il un $A$-Module projectif ? (Je me rappelle avoir posé la question à D. Lazard, mais je ne me rappelle pas s'il l'évait résolue dans un sens ou dans l'autre. Je crois que non). Il faudrait donc à la prochaine occasion faire un errata pour 17.1.6 , dans lequel on \deleted{soulèverait} signalerait également la formulation équivalente que je viens de signaler pour la question du caractère local de la lissité formelle. - Si la réponse devait être négative, il y aurait lieu de considérer aussi la notion de morphisme localement formellement lisse. En fait, dans tous les cas que j'ai rencontrés, \deleted{quand} on prouve toujours directement la lissité formelle globale, telle qu'elle est définie dans le texte.

\deleted{J'ai fait lire mon projet de par.0 pour la réédition de EGA I à Samuel et Illusie, dont je te transmets certaines critiques. Page 4, au lieu de ``quasi-totalité'', il est plus prudent de dire ``plupart''. Page 8, j\deleted{e}'utilise tantôt la terminologie ``idéal racine'', tantôt ``idéal radical'' et j'écris $\operatorname{rad}(J)$, - il faut unifier. Page 9, ligne 4, style vicieux, lire ``comme anneaux de valeurs pour les coordonnées de solutions ... ''. Page 10, donner en note de bas de page le sens du mot}
% archive sheet number: 93
